%======================================================================
% CardamomOT mixed-effects networks: project note
%======================================================================
\documentclass[a4paper, 11pt]{article}
\usepackage[utf8]{inputenc}
\usepackage[T1]{fontenc}
\usepackage{amsmath, amssymb, amsthm}
\usepackage[margin=1in]{geometry}
\usepackage{enumitem}
\usepackage{booktabs}
\usepackage{hyperref}
\hypersetup{colorlinks=true, linkcolor=blue, citecolor=blue, urlcolor=blue}

\DeclareMathOperator*{\argmin}{arg\,min}
\DeclareMathOperator{\diag}{diag}
\DeclareMathOperator{\tr}{tr}
\newcommand{\E}{\mathbb{E}}
\newcommand{\pen}{\mathrm{pen}}
\newcommand{\tb}{\bar\theta}
\newcommand{\tofill}[1]{\par\smallskip\noindent\emph{[To be completed: #1]}\par\smallskip}

\theoremstyle{plain}
\newtheorem{prop}{Proposition}
\theoremstyle{remark}
\newtheorem{rem}{Remark}

\title{CardamomOT-ME\\[0.3em] \large Gene regulatory networks across patients:\\ calibrated network conditions and a mixed-effects model}
\author{Project note}
\date{\today}

\begin{document}
\maketitle

%======================================================================
\section*{Project note}
%======================================================================

\paragraph{Objective.} CardamomOT~\cite{mauge2026cardamomot} infers a gene regulatory network from single-cell time series with an EM-type loop that alternates between trajectories and modes on one side and a sparse regression of each target gene on the other. With several patients, one wants a population network, the network of each patient and a measure of how much patients differ, without tuning a coupling penalty by hand.

\paragraph{The thread.} Two steps. First, the penalties of the existing network conditions are made the direct generalisation of the single-network method: every level is a quantile universal threshold (QUT)~\cite{giacobino2017qut}, so that a false edge is as unlikely per condition as in a network inferred alone. Second, for patients, the deviations from the mean network become a random effect: $\psi_p = \tb + \varepsilon_p$, $\varepsilon_p \sim \mathcal N(0, \tau^2)$, with $\tau^2$ learned in the loop.

\paragraph{Why a random effect.} A patient is a draw from a population: the targets are $\tb$ and $\tau^2$, and each patient has one vote whatever its number of cells or time points (Proposition~\ref{prop:mstep}). A treatment or a cell line is a fixed effect, for which Section~\ref{sec:fixed} remains the right model. The coupling is learned instead of set by hand, shrinkage follows the information of each patient on each edge, and a new patient can be simulated from $\mathcal N(\tb, \tau^2)$.

\paragraph{Where we stand.} Both steps are implemented in the package (option \texttt{mixed\_effect\_model\_conditions}). Given the trajectories and the modes, they behave as designed through the real network fit: false edges controlled, $\hat\tau$ calibrated. Over the full loop on FN8, $\hat\tau$ follows the true heterogeneity once the mixtures of the samples are harmonised, but the mixed model does not improve on the fixed-effect model with $4$ patients, and no method recovers the deviations of the patients (Section~\ref{sec:pipeline}). The largest source of differences between replicates is the mixture of each sample, not the network (Section~\ref{sec:open}).

%======================================================================
\section{The network fit}\label{sec:fit}
%======================================================================

At each update of the loop, the network is fitted target by target. For a target gene, with rows $i$ (trajectory states), regulators $Y_i \in \mathbb R^K$, sigmoid targets $y_i \in [0,1]$ and per-sample basals $b_s$,
\begin{equation}\label{eq:data}
L(\theta, b) = \sum_{s} w_s \sum_{i \in s} \mathrm{CE}\big(\sigma(b_s + Y_i^\top \theta),\, y_i\big),
\end{equation}
plus a flow-matching term of the same form along the intervals. With one network, the sparsity level is the QUT of the target: the level-quantile of $\max_j |\partial_j L|$ at the null model (no interaction, basals fitted), the target being permuted within each sample, which is the zero-thresholding level of the L1 penalty under the null. A false edge then enters with probability about $1 -$ \texttt{qut\_level} per target. The sample weights $w_s$ are set by \texttt{alpha} ($\alpha = 1$: each sample has the same total mass).

%======================================================================
\section{Network conditions as fixed effects}\label{sec:fixed}
%======================================================================

With conditions $c = 1, \dots, C$, the network of condition $c$ is $\theta_c = \theta_s + \Delta_c$. Write $g_c \in \mathbb R^K$ for the gradient of the data term of condition $c$ at the null, on the scale of the joint objective. We use
\begin{equation}\label{eq:pen}
s\,\|\theta_s\|_1 + \sum_c a_c \|\theta_c\|_1 + \sum_c f_c \|\Delta_c\|_1,
\end{equation}
\begin{equation}\label{eq:levels}
s = \kappa\,\rho\,\lambda, \qquad a_c = \kappa\,(1 - \rho)\,\lambda_c, \qquad f_c = \kappa\,\pen\,\lambda_c, \qquad \rho = \min(\pen, 1),
\end{equation}
where $\lambda$ is the QUT of $\sum_c g_c$ (the level of a common network), $\lambda_c$ the QUT of $g_c$ with the same permutations (the level of condition $c$ inferred alone), and $\kappa \ge 1$ is defined below. The parameter $\pen$ is \texttt{network\_condition\_pen}.

\begin{prop}[Routes of a false edge]\label{prop:routes}
For one edge with gradients $(g_c)$ at zero, the zero solution of~\eqref{eq:pen} is optimal if and only if there is $u \in [-1,1]^C$ with $|g_c - f_c u_c| \le a_c$ for all $c$ and $|\sum_c f_c u_c| \le s$. In particular, for every subset $A$ of conditions, it requires
\begin{equation}\label{eq:routes}
\Big|\sum_{c \in A} g_c\Big| \le \sum_{c \in A} a_c + \min\Big(\sum_{c \in A} f_c,\; s + \sum_{c \notin A} f_c\Big).
\end{equation}
\end{prop}
\begin{proof}[Sketch]
The conditions are the stationarity conditions in $\theta_c$ and $\theta_s$ with subgradients $u_c$ of $|\theta_c - \theta_s|$. For~\eqref{eq:routes}, sum over $A$ the first conditions and bound $\sum_A f_c u_c$ in two ways, directly and through $\sum_c f_c u_c$.
\end{proof}

Three consequences motivate~\eqref{eq:levels}.
\begin{itemize}[leftmargin=*]
\item \emph{Each level is calibrated on its own gradient.} The noise of $g_c$ is of order $\lambda / \sqrt C$ for balanced conditions, not $\lambda / C$. A penalty $\omega_c \lambda$ on $\theta_c$ with $\sum_c \omega_c = 1$, which gives the right cost to a common edge, is therefore too low by $\sqrt C$ for a condition-specific one: for $\pen$ below about $1/\sqrt C - 1/C \approx 0.24$ (for $C = 3$ and $6$) the networks are denser than independent ones.
\item \emph{A direct penalty on the shared network closes the indirect routes.} Without it ($s = 0$), a small condition with a large $\lambda_c$ can bypass its level: take $A = \{c\}$ in~\eqref{eq:routes}, whose right-hand side is at most $a_c + \sum_{c' \ne c} f_{c'}$. With $s = \lambda$, which is larger than $\lambda_c$ as the noise of a sum of independent gradients, the minimum in~\eqref{eq:routes} for $A = \{c\}$ is $f_c$, and no indirect route is cheaper than the direct one.
\item \emph{Both limits hold.} For $\pen = 0$ each network is at its own QUT and the fit reduces to independent fits. For $\pen \to \infty$ only the shared network is left, at the QUT of a common network. With $C = 1$ the method is the single-network one. At $\pen = 1$ the shared network and the deviations are each at their own QUT, which is the default. The form at $\pen = 1$ is the dirty model of multi-task regression~\cite{jalali2010dirty}, close to the fused penalties of joint graphical models~\cite{danaher2014joint}.
\end{itemize}

\paragraph{The scale $\kappa$.} At $\pen = 1$, a false edge has two main routes, and two families of tests at level \texttt{qut\_level} give more false edges than one. The QUT is the quantile of the zero-thresholding function, here the smallest common scale of the penalties that keeps the network of condition $c$ at zero through the two main routes of~\eqref{eq:routes}:
\begin{equation}
\kappa_c = \mathrm{quantile}\Big(\max\Big(\frac{\|\sum_{c'} g_{c'}\|_\infty}{s_0 + \sum_{c'} a_{c',0}},\ \Big\|\frac{g_c}{a_{c,0} + f_{c,0}}\Big\|_\infty\Big)\Big), \qquad \kappa = \frac 1 C \sum_c \kappa_c,
\end{equation}
where the index $0$ denotes the levels before scaling. By construction $\kappa = 1$ in both limits. Implementation: \texttt{qut\_conditions} in \texttt{inference/network.py}. The refinement (garrote) keeps one correction per edge, shared by the conditions.

%======================================================================
\section{The mixed model}\label{sec:mixed}
%======================================================================

\subsection{Model}

One network condition is one patient, $p = 1, \dots, P$. For each target gene $g$, that is one column of the network, regressed on its regulators,
\begin{equation}\label{eq:model}
\psi_p = \tb + \varepsilon_p, \qquad \varepsilon_{pj} \sim \mathcal N(0, \tau_g^2),
\end{equation}
in the units of the effective network, the proteins of all genes living on one scale. The objective at fixed $\tau_g^2$ is
\begin{equation}\label{eq:objective}
L(\psi, b) + \lambda_m \|\tb\|_1 + \frac 1 {2\tau_g^2} \sum_p \|\varepsilon_p\|^2 + \frac 1 {2\tau_b^2} \sum_p (\bar b_p - \bar b)^2 ,
\end{equation}
the last term being the random effect on the basals (Section~\ref{sec:basal}), with the ridge of \texttt{network\_ridge} on $\tb$.

\subsection{One vote per patient}

\begin{prop}\label{prop:mstep}
In the complete-data likelihood of~\eqref{eq:model}, the M-step of the mean network is the unweighted mean $\tb = \frac 1 P \sum_p \E[\psi_p \mid \text{data}]$. Under a quadratic approximation of $L$ with curvature $H_p$ at the estimate $\hat\psi_p$ of patient $p$ alone, its fixed point is the precision-weighted mean of the $\hat\psi_p$ with weights $(H_p^{-1} + \tau^2 I)^{-1}$.
\end{prop}

As in random-effects meta-analysis, the weights do not depend on the number of cells once $H_p^{-1} \ll \tau^2 D^{-1}$: a patient with $5000$ cells weighs as much as a patient with $500$. A patient with too few cells to say anything weighs less, on the edges it says nothing about. So $\alpha$ is no longer needed to balance patients, and it is ignored by the mixed model.

The remaining question is how much information a cell brings, which sets $H_p$. Taken as a tempered likelihood~\cite{holmes2017power}, $\alpha = 1$ means that every sample has the same effective size. It breaks when it inflates a small sample: a sample of $50$ cells among samples of $50$ to $5000$ counts as about $1300$ cells. We use the cap
\begin{equation}\label{eq:cap}
c_s = \min(1, n_{\mathrm{ref}} / n_s), \qquad n_{\mathrm{ref}} = \mathrm{median}_s\, n_s,
\end{equation}
which never inflates a sample and states that cells beyond $n_{\mathrm{ref}}$ bring no further information. Table~\ref{tab:cells} compares these choices.

\begin{table}[h]
\centering\small
\begin{tabular}{lcccc}
\toprule
 & $50$--$5000$ cells & $300$--$5000$ cells & $50$--$5000$, at most $200$ indep.\ cells & $50$--$5000$, $\tau = 0$ \\
\midrule
Common, $\alpha = 0$ & -- / 0.312 & -- / 0.276 & -- / 0.362 & -- / 0.041 \\
Common, $\alpha = 1$ & -- / 0.265 & -- / 0.259 & -- / 0.271 & -- / 0.113 \\
Mixed, $\alpha = 0$ & 0.49 / 0.221 & 0.50 / 0.208 & 0.57 / 0.262 & 0.05 / 0.043 \\
Mixed, $\alpha = 1$ & 0.79 / 0.343 & 0.52 / 0.212 & 0.84 / 0.392 & 0.36 / 0.184 \\
Mixed, cap~\eqref{eq:cap} & 0.46 / 0.218 & 0.49 / 0.207 & 0.52 / 0.243 & 0.06 / 0.068 \\
\bottomrule
\end{tabular}
\caption{Toy logistic mixed model, $6$ patients, $8$ regulators, $30$ targets, $4$ seeds, true $\tau = 0.5$ unless stated: $\hat\tau$ / error on $\tb$ (best possible about $0.20$ for $\tau = 0.5$). With a common network, $\alpha = 1$ is better as soon as patients differ, which is why it is kept for the fixed-effect path.}\label{tab:cells}
\end{table}

\subsection{Learning the variances}

$\tau^2$ cannot be optimised jointly with $\varepsilon$: the joint maximum is degenerate ($\varepsilon \to 0$, $\tau \to 0$). The deviations are latent variables, next to the modes of the EM. With a Laplace approximation~\cite{laird1982random, breslow1993approximate}, the fit~\eqref{eq:objective} at $\tau^2$ gives $\hat\varepsilon_p$ and the joint Gauss-Newton matrix $H$ of the kept entries, and the REML score equation of $\tau^2$ is the fixed point of
\begin{equation}\label{eq:tau}
T(\tau^2) = \frac{\sum_p \sum_{j} \big(\hat\varepsilon_{pj}^2 + V_{pj}\big)}{n}, \qquad V = H^{-1},
\end{equation}
over the $n$ entries of the edges of the support of $\tb$. Taking $V_p$ as the block of the inverse of the \emph{joint} matrix integrates out $\tb$ and the basals: this is REML~\cite{patterson1971recovery}, which removes the bias of about $(P-1)/P$ of maximum likelihood, large with few patients (Table~\ref{tab:saem}). The term in $V$ plays for $\tau^2$ the role that the posterior variance plays for the precision of the network prior on the modes.

\paragraph{Schedule.} The variances are not updated along the loop. \emph{Burn-in}: the fixed-effect model of Section~\ref{sec:fixed} at \texttt{network\_condition\_pen}, which needs no variance, until \texttt{min\_n\_loops}; independent networks there inflate $\hat\tau$, each patient having shaped its own trajectories. \emph{Variances}: learned once, on the trajectories of that iteration, frozen, with all the states. \emph{Rest of the loop}: the fit~\eqref{eq:objective} at these fixed values.

\paragraph{Solver.} Iterating $\tau^2 \leftarrow T(\tau^2)$ is an EM, too slow near $0$, where its limit depends on the number of steps. The score equation $T(\tau^2) = \tau^2$ is solved by bisection on $\log \tau^2$ in $[10^{-4}, 10^2]$, $7$ fits in a pass, first for the pooled $\tau_0^2$ and $\tau_b^2$, then for each target gene with the shrinkage
\begin{equation}\label{eq:shrink}
\tau_g^2 = \frac{s_g(\tau_g^2) + d_0\, \tau_0^2}{n_g(\tau_g^2) + d_0}, \qquad d_0 = P / 2,
\end{equation}
that is an inverse-gamma prior worth half an edge. A target without edge gets $\tau_0^2$.

\paragraph{Why not SAEM.} SAEM~\cite{kuhn2004coupling} is the classical algorithm for these models. In our setting it brings nothing on $\tau$ from about $100$ cells per patient: the posterior of $\varepsilon_p$ is log-concave given the modes and cells are numerous. It needs several chains with few patients: with one chain and $\gamma = 1$, $\tb$ is the mean of the draws, which loses a factor $(P-1)/P$ per iteration when the data are weak, and $\hat\tau$ collapses. A draw in the network term would also bring back noise in the basins, which were made deterministic. Sampling would become useful for non-Gaussian deviations~\cite{park2008bayesian}, at the kink of which the Laplace approximation fails, or with very few cells per patient, where Laplace maximum likelihood is biased downwards~\cite{breslow1995bias} and where nothing was reliable in our tests.

\begin{table}[h]
\centering\small
\begin{tabular}{lccc}
\toprule
Cells per patient & Laplace ML & SAEM ML & Laplace REML \\
\midrule
$50$ to $5000$ & 0.442 & 0.447 & 0.503 \\
$300$ to $5000$ & 0.460 & 0.462 & 0.509 \\
$100$ to $5000$, steep sigmoids & 0.442 & 0.444 & 0.503 \\
$200$ to $5000$, $3$ patients & 0.390 & 0.389 & 0.495 \\
$100$ for all & 0.405 & 0.464 (10 chains) & 0.488 \\
$30$ for all & 0.329 & 0.49 or 0.07 depending on the initial value (10 chains) & 0.500 \\
\bottomrule
\end{tabular}
\caption{Toy logistic mixed model, one seed, true $\tau = 0.5$: $\hat\tau$. SAEM with one chain and $300$ iterations in the first four rows.}\label{tab:saem}
\end{table}

\subsection{A QUT for the mean network}

The objective~\eqref{eq:objective} changes the zero-thresholding function, and the single-network QUT cannot be used: it kills the mean network.

\begin{prop}\label{prop:qut}
With $s_p$ the score of patient $p$ at the null, $H_p$ its Gauss-Newton matrix with the basals eliminated, and $P_\tau = I / \tau^2$, a quadratic approximation gives, at $\tb_j = 0$ with the deviations fitted,
\begin{equation}
\partial_{\tb} = \sum_p (I + H_p P_\tau^{-1})^{-1} s_p .
\end{equation}
For $\tau \to 0$ it is the common gradient $\sum_p s_p$, and for $H_p P_\tau^{-1} \gg I$ it is $\tau^{-2} \sum_p \hat\psi_p$.
\end{prop}

$\lambda_m$ is the level-quantile of the max norm of this statistic under the null of the model, the edge being absent on average but varying between patients: $s_p$ is the score of the target permuted within each sample plus $H_p e_p$, with $e_p \sim \mathcal N(0, P_\tau^{-1})$. Permutations alone (no edge in any patient) ignore the variance between patients and let false population edges enter (Table~\ref{tab:qut}). The threshold on $\tb_j$ is about $z\, \tau / \sqrt P$: the cost of few patients is a loss of power on weak edges.

\subsection{Sparsity of the patient networks}

The deviations are Gaussian, so they are not sparse. The networks of the patients are, through the following rule.
\begin{enumerate}[leftmargin=*]
\item On the support of $\tb$, $\varepsilon_p$ is kept, shrunk.
\item Out of it, $\varepsilon_{pj} = 0$ unless its posterior $z_{pj} = \hat\varepsilon_{pj} / \sqrt{V_{pj}}$ exceeds the universal threshold at \texttt{qut\_level} for the $K P$ tests of the target: these are the patient-specific edges.
\item The network is refitted on the kept entries, with its L1.
\item For the statistics of $\tau^2$ only, the same entries are refitted with the L1 relaxed on the support of $\tb$: with the L1, $\tb$ is pulled towards $0$ and the deviations of all patients the other way, which inflates $\tau^2$ through a self-reinforcing loop. Relaxed in the network itself, it misleads the refinement, whose level follows the magnitude of the kept edges: over the loop at $\tau \to 0$, an AUPR of $0.31$ against $0.60$ with the L1 ($0.63$ for a common network).
\item $\tau^2$ is estimated on the support of $\tb$ only, the patient-specific edges being selected, hence biased.
\end{enumerate}
The posterior $z$ of each edge in each patient is saved (\texttt{network\_deviation\_z.npy}) and measures the differences between patients.

\subsection{Basals}\label{sec:basal}

Basals stay per sample (\texttt{dataset\_id}). A free fixed effect per sample nested in a random patient effect is not identifiable: adding a constant to the patient effect and subtracting it from the fixed effects of its samples leaves the likelihood unchanged, and the prior then sets the patient effect, and $\tau_b$, to $0$. The identifiable form is the within-between decomposition~\cite{mundlak1978pooling}: $b_s = \bar b + v_p + f_s$ with $\sum_{s \in p} f_s = 0$, in which only the patient means are penalised, as in~\eqref{eq:objective}. Contrasts between the samples of a patient remain free, a patient with one sample is shrunk as a whole, and no new parameter is needed. $\tau_b^2$ follows the same score equation, with $P - 1$ degrees of freedom. As the regulators are not centred, a deviation of a patient also moves its basal, so $\tau_b$ grows with $\tau$. Samples pinned by a KO/OV prior and single-mode samples of the target are left out.

\begin{rem}[Cell types]
A patient is exogenous: it does not change along a trajectory. A cell type is endogenous: the types are the attractors of the network, and the change of burst frequency between types is what the network must explain through its regulators. A per-type effect would absorb the regulation, and the simulation would become a mixture of experts. It makes sense only for labels fixed along a trajectory (separated lineages, clone, genotype), which are then crossed effects with a second variance component.
\end{rem}

%======================================================================
\section{Validation}\label{sec:validation}
%======================================================================

\subsection{Independent networks at $\pen = 0$}

$\pen = 0$ goes through one single-network fit per condition, bit-identical to the single-network path at the level of the network fit. Over \texttt{fit\_network}, we compared two FN8 datasets inferred jointly to each inferred alone, with the difference between two seeds of a run alone as reference, the second dataset having its genes permuted so that the two networks really differ ($9$ common edges out of $17$).

\begin{center}\small
\begin{tabular}{lcc}
\toprule
 & Seed noise (alone vs alone) & Joint vs alone \\
\midrule
Condition A, same support & 0.844 & 0.81--0.89 \\
Condition A, correlation of the values & 0.976 & 0.85--0.96 \\
Condition A, AUPR & 0.81 / 0.77 (alone) & 0.81 / 0.78 (joint) \\
Condition B, same support & 0.891 & 0.89--0.94 \\
Condition B, correlation of the values & 0.989 & 0.93--0.98 \\
Condition B, AUPR & 0.82 / 0.81 (alone) & 0.87 / 0.86 (joint) \\
\bottomrule
\end{tabular}
\end{center}

The structure is equivalent up to the seed noise, and values differ slightly more. What remains shared in the loop at $\pen = 0$: the family of $\sigma$ of each target (\texttt{model\_learning\_grn}), the thresholds of the regulators (mean over the samples), the stopping of the loop.

\subsection{Network conditions}

Through the real \texttt{main\_loop\_inference} (final networks), $6$ conditions without real differences except where stated, $12$ regulators, $24$ targets, one seed; $500$ cells per condition, or $150$ to $3000$ in the unbalanced row.

\begin{center}\small
\begin{tabular}{lcccc}
\toprule
 & $\pen = 0.1$ & $\pen = 1$ & Independent & Common \\
\midrule
No signal: false edges per condition & 0.53 & 0.54 & 0.83 & 0.71 \\
No signal, unbalanced: worst condition & 0.67 & 0.62 & 1.29 & 0.79 \\
Weak common edges: recall & 0.76 & 0.99 & 0.65 & 1.00 \\
Condition-specific edges: recall & 0.98 & 0.98 & 1.00 & 0.31 \\
Condition-specific edges: false edges per condition & 0.17 & 0.38 & 0.00 & 1.57 \\
\bottomrule
\end{tabular}
\end{center}

Over three seeds, the results are stable. As $\pen \to 0$, $99\,\%$ of the support is that of independent fits; the difference comes from the shared garrote, on the conservative side (false edges $0.53$ vs $0.83$, recall $0.98$ vs $1$). For $\pen = 10^4$ all conditions are exactly equal, with a support $92$ to $100\,\%$ that of the common fit.

Cost of a value above $1$, same scenarios, three seeds:

\begin{center}\small
\begin{tabular}{lccc}
\toprule
 & $\pen = 1$ & $\pen = 2$ & $\pen = 3$ \\
\midrule
No signal: false edges per condition & 0.45 & 0.51 & 0.50 \\
No signal: false deviations per condition & 0.25 & 0.00 & 0.00 \\
No signal, unbalanced: worst condition & 0.61 & 0.57 & 0.58 \\
Weak common edges: recall & 1.00 & 1.00 & 1.00 \\
Condition-specific edges: recall & 0.96 & 0.85 & 0.60 \\
Condition-specific edges: false edges per condition & 0.32 & 1.52 & 2.01 \\
\bottomrule
\end{tabular}
\end{center}

Above $\pen = 1$, an edge specific to one condition that cannot pay for its deviation goes through the shared network and appears in the other conditions: this is the price of fewer false differences.

\subsection{Mixed model}

Through the real \texttt{main\_loop\_inference}, with the solver of Section~\ref{sec:mixed}, $12$ regulators of which $3$ true edges, $24$ targets, Bernoulli targets, two seeds, regulators of standard deviation $0.2$.

\begin{table}[h]
\centering\small
\begin{tabular}{llcccccc}
\toprule
Design & Scenario & True $\tau$ & $\hat\tau_0$ & False edges of $\tb$ & Recall of $\tb$ & Patient edges & Recall \\
\midrule
$6$ patients & Replicates & 0 & 0.01 & 0.00 & 1.00 & 0.00 & -- \\
$300$--$2000$ & Deviations on true edges & 1.25 & 1.16 & 0.00 & 1.00 & 0.00 & -- \\
cells & & 2.5 & 2.40 & 0.00 & 0.97 & 0.33 & -- \\
 & & 5 & 4.83 & 0.00 & 0.94 & 0.67 & -- \\
 & + 2 patient-specific edges & 2.5 & 2.53 & 0.04 & 0.99 & 2.04 & 0.96 \\
 & No signal & 0 & 0.05 & 0.46 & -- & 0.00 & -- \\
\midrule
$3$ patients & Replicates & 0 & 0.01 & 0.00 & 1.00 & 0.00 & -- \\
$500$--$1500$ & Deviations on true edges & 1.25 & 1.04 & 0.00 & 1.00 & 0.00 & -- \\
cells & & 2.5 & 2.21 & 0.00 & 0.95 & 0.10 & -- \\
 & & 5 & 4.33 & 0.00 & 0.95 & 0.17 & -- \\
 & + 2 patient-specific edges & 2.5 & 2.67 & 0.40 & 0.99 & 1.56 & 0.76 \\
 & No signal & 0 & 0.01 & 0.65 & -- & 0.00 & -- \\
\bottomrule
\end{tabular}
\caption{Mixed model through the real fit. Per target: false edges of the mean network and patient-specific edges. Without signal, the false edges are those of the level of the QUT.}\label{tab:qut}
\end{table}

In the toy model with $6$ patients, permutations alone give $1.35$ to $6.60$ false edges of $\tb$ per target against $0$ to $0.60$ with the QUT of the model, and reusing the single-network QUT gives a recall of $0$.

\paragraph{Calibration of $\tau$.} With the bisection, replicates reach the bottom of the bracket at once ($\hat\tau$ of $0.01$ to $0.02$), where EM steps were still decreasing. Soft targets equal to the expected values with noise of sd $0.1$: $\hat\tau$ of $2.04$ for $2.5$, and $0.01$ for $0$; sd $0.3$: $1.25$ and $0.14$. In the toy model, cells replicated $5$ to $25$ times lead to $\hat\tau$ of $0.58$ instead of $0.48$: the absolute scale of the likelihood matters, moderately. $\tau_b$ of the replicates is $0.50$ to $0.53$ for a true $0.5$, and grows with $\tau$ as the regulators are not centred ($4.1$ for $\tau = 5$, against $\sqrt{0.25 + 0.75\,\tau^2} = 4.4$ expected in this design).

\paragraph{One $\tau$ per target.} Half of the targets at $\tau = 1.25$, half at $5$: root mean square error on $\log\tau$ without shrinkage / with~\eqref{eq:shrink}, $0.22$ / $0.32$ with $6$ patients and $3$ edges per target, $0.50$ / $0.41$ with $3$ patients. With $4$ patients and one edge per target, the same error goes from $1.17$ to $0.62$, and a target of the replicates reaches $\hat\tau = 10$ without shrinkage, none with it. Larger priors ($d_0 \ge 2P$) pull the low targets too far up, the prior being centred on an arithmetic mean of variances. A prior on $\log\tau^2$ gives the same results.

\subsection{Full loop}\label{sec:pipeline}

FN8 simulated by the model of the model-learning benchmark (proteins hidden, mixed families of $\sigma$, single-sample AUPR of about $0.6$ for this run), $4$ patients of $400$ to $1200$ cells, with $\theta_{ij} (1 + \rho Z)$ on the true edges and basals shifted by $\mathcal N(0, \rho^2)$. One seed, \texttt{infer\_network\_structure} only, with per-sample mixtures (\texttt{integrate\_samples} off) or harmonised ones (on). AUPR of the mean network / per patient, kept edges and false edges per patient ($17$ true edges), correlation between the inferred and true deviations on the true edges.

\begin{table}[h]
\centering\small
\begin{tabular}{lllcccc}
\toprule
Mixtures & Data & Method & AUPR mean / patient & Edges (false) & Corr.\ deviations & Edges differing \\
\midrule
Per sample & $\rho = 0$ & Mixed ($\hat\tau = 1.32$) & 0.46 / 0.46 & 11.0 (7.0) & -- & 11 \\
 & & $\pen = 1$ & 0.51 / 0.46 & 25.8 (19.5) & -- & 25 \\
 & & $\pen = 2$ & 0.46 / 0.47 & 27.2 (21.2) & -- & 5 \\
 & & $\pen = 3$ & 0.38 / 0.38 & 27.0 (21.0) & -- & 0 \\
 & & Independent ($\pen = 0$) & 0.55 / 0.46 & 23.8 (16.8) & -- & 39 \\
 & & Common ($\pen = 1000$) & 0.44 / 0.44 & 25.0 (18.0) & -- & 0 \\
\cmidrule{2-7}
 & $\rho = 0.3$ & Mixed ($\hat\tau = 2.04$) & 0.33 / 0.34 & 8.2 (6.2) & $-0.12$ & 9 \\
 & & $\pen = 1$ & 0.43 / 0.40 & 29.0 (22.0) & 0.12 & 27 \\
 & & $\pen = 2$ & 0.42 / 0.42 & 24.8 (19.2) & $-0.29$ & 8 \\
 & & $\pen = 3$ & 0.43 / 0.43 & 26.0 (20.0) & -- & 0 \\
 & & Independent ($\pen = 0$) & 0.72 / 0.55 & 28.0 (19.2) & 0.01 & 47 \\
 & & Common ($\pen = 1000$) & 0.40 / 0.40 & 25.0 (19.0) & -- & 0 \\
\cmidrule{2-7}
 & $\rho = 0.6$ & Mixed ($\hat\tau = 2.27$) & 0.47 / 0.47 & 7.0 (4.0) & 0.30 & 7 \\
 & & $\pen = 1$ & 0.44 / 0.45 & 23.0 (16.2) & $-0.20$ & 25 \\
 & & Independent ($\pen = 0$) & 0.42 / 0.42 & 25.8 (17.8) & 0.10 & 47 \\
 & & Common ($\pen = 1000$) & 0.41 / 0.41 & 22.0 (16.0) & -- & 0 \\
\midrule
Harmonised & $\rho = 0$ & Mixed ($\hat\tau = 0.32$) & 0.58 / 0.57 & 32.0 (24.0) & -- & 30 \\
 & & $\pen = 1$ & 0.62 / 0.62 & 25.5 (18.2) & -- & 10 \\
 & & Independent ($\pen = 0$) & 0.61 / 0.54 & 23.5 (14.5) & -- & 36 \\
 & & Common ($\pen = 1000$) & 0.63 / 0.63 & 34.0 (26.0) & -- & 0 \\
\cmidrule{2-7}
 & $\rho = 0.3$ & Mixed ($\hat\tau = 0.56$) & 0.56 / 0.55 & 22.0 (16.0) & 0.17 & 22 \\
 & & $\pen = 1$ & 0.61 / 0.61 & 27.2 (19.2) & 0.07 & 12 \\
 & & Independent ($\pen = 0$) & 0.73 / 0.61 & 29.0 (20.2) & $-0.01$ & 43 \\
 & & Common ($\pen = 1000$) & 0.63 / 0.63 & 24.0 (15.0) & -- & 0 \\
\cmidrule{2-7}
 & $\rho = 0.6$ & Mixed ($\hat\tau = 1.06$) & 0.39 / 0.41 & 9.5 (6.5) & 0.08 & 11 \\
 & & $\pen = 1$ & 0.51 / 0.49 & 23.5 (16.8) & 0.30 & 18 \\
 & & Independent ($\pen = 0$) & 0.70 / 0.52 & 18.5 (12.0) & 0.24 & 39 \\
 & & Common ($\pen = 1000$) & 0.52 / 0.52 & 20.0 (13.0) & -- & 0 \\
\bottomrule
\end{tabular}
\caption{FN8 over the full loop.}\label{tab:pipeline}
\end{table}

\begin{itemize}[leftmargin=*]
\item \emph{The mixtures are the largest source of differences between replicates.} With per-sample mixtures, the same gene gets $2$ to $5$ levels and a mean target of $0.49$ to $0.83$ across replicates of one network. Harmonising them raises the AUPR per patient of the common network from $0.44$ to $0.63$ and of $\pen = 1$ from $0.46$ to $0.62$, and lowers the false differences of $\pen = 1$ from $25$ to $10$. They may also carry real differences between patients, which a projection on a reference sample removes.
\item \emph{$\hat\tau$ follows $\rho$ with harmonised mixtures}: $0.32$, $0.56$, $1.06$. Its floor on replicates is above $0$, and its absolute value is hard to compare to the truth, the inferred network being on another scale (median $|\theta|$ of about $3.5$ on the true edges, $10$ in the simulator). With per-sample mixtures it is high, and unstable: the score is close to its fixed point from $\tau = 1$ to $5$, and a small change of the run moved $\hat\tau$ from $0.86$ to $1.32$ on replicates and from $7.0$ to $2.0$ at $\rho = 0.3$.
\item \emph{The variance comes from the patient level.} On frozen trajectories with per-sample mixtures, $T(\tau^2) > \tau^2$ up to $\tau \approx 1$ on replicates, and $T(\tau^2) < \tau^2$ everywhere with random pseudo-patients, whether the label is drawn per state or per trajectory: the correlation of the states along a trajectory is not the cause. It is not an imprint of the patient networks either: the trajectories of a common network give about the same $\hat\tau$. Common amplitudes and $\tau_b$ do not change it. With harmonised mixtures and the trajectories of a common network, the excess nearly vanishes on replicates (at most $1\,\%$ above the fixed point, up to $\tau = 0.2$) and remains at $\rho = 0.3$ ($6\,\%$, up to $\tau = 0.7$).
\item \emph{No gain of the mixed model with $4$ patients.} Its AUPR stays below that of $\pen = 1$ and of the common network (harmonised: $0.57$, $0.55$, $0.41$ per patient against $0.62$, $0.61$, $0.49$). The threshold on $\tb$ is about $z\,\tau / \sqrt P$: at $\rho = 0.6$ only $9.5$ edges are kept.
\item \emph{The deviations are not recovered by any method.} Correlations stay at or below $0.30$, one seed. No entry of the posterior $z$ exceeds $3$, even at $\rho = 0.6$.
\item \emph{The mean of independent networks stays the best mean network when patients differ} ($0.73$ and $0.70$ for $\rho = 0.3$ and $0.6$), though not the best network of each patient.
\item \emph{Noise between patients in the fixed-effect model.} The QUT of the deviations only knows the noise of the cells: $25$ false differences at $\pen = 1$ on replicates with per-sample mixtures, $5$ at $\pen = 2$, $10$ at $\pen = 1$ with harmonised mixtures. Values above $1$ have a cost on the edges specific to a condition (Section~\ref{sec:validation}), so the default stays at $1$.
\end{itemize}

%======================================================================
\section{Open points}\label{sec:open}
%======================================================================

\begin{itemize}[leftmargin=*]
\item \emph{Mixtures of the patients (main point).} Per-sample grids make replicates differ. Harmonising only the number of levels between patients, or their selection, could keep real differences of level values while removing most of this noise.
\item \emph{More patients.} The mixed model has no reason to win with $4$ patients: its threshold falls as $1 / \sqrt P$. FN8 with about $10$ patients and harmonised mixtures is the next test.
\item \emph{A benchmark where deviations are identifiable.} Observed proteins, to separate what the method cannot do from what the data cannot show.
\item \emph{Floor of $\hat\tau$ on replicates.} A burn-in with a common network gives a flatter score on replicates; to be tested through the loop.
\item \emph{Coupling at $\pen = 0$.} Family of $\sigma$ and thresholds of the regulators per condition would need one specification per condition everywhere in the package.
\item \emph{Calibration on real data.} The loss is a pseudo-likelihood (soft targets, sub-samples, cells correlated along trajectories). Splitting one sample into pseudo-patients is the check to run on real data: $\hat\tau$ must be close to $0$. The Gauss-Newton matrix with the mRNA memory (\texttt{MemStates}) has not been run yet.
\item \emph{Not yet supported.} \texttt{select\_genes\_per\_condition} and \texttt{hard\_forcing\_ref} with the mixed model; fixed effects crossed with patients (treatment within patient).
\item \emph{Sparse deviations.} The hard rule on $z$ could become a two-group model with a learned proportion of deviating edges.
\item \emph{New patients.} Simulating from $\mathcal N(\tb, \tau^2)$ and validating on a patient left out.
\end{itemize}

\begin{thebibliography}{99}

\bibitem{mauge2026cardamomot}
Y.~Maugé, E.~Ventre.
\newblock CardamomOT: a mechanistic optimal transport-based framework for gene regulatory network inference, trajectory reconstruction and generative modeling.
\newblock \emph{bioRxiv}, 2026. doi:10.64898/2026.03.31.715390. Accepted in \emph{PLoS Computational Biology}.

\bibitem{giacobino2017qut}
C.~Giacobino, S.~Sardy, J.~Diaz-Rodriguez, N.~Hengartner.
\newblock Quantile universal threshold.
\newblock \emph{Electronic Journal of Statistics}, 11(2):4701--4722, 2017.

\bibitem{jalali2010dirty}
A.~Jalali, P.~Ravikumar, S.~Sanghavi, C.~Ruan.
\newblock A dirty model for multi-task learning.
\newblock \emph{Advances in Neural Information Processing Systems}, 23, 2010.

\bibitem{danaher2014joint}
P.~Danaher, P.~Wang, D.~M.~Witten.
\newblock The joint graphical lasso for inverse covariance estimation across multiple classes.
\newblock \emph{Journal of the Royal Statistical Society: Series B}, 76(2):373--397, 2014.

\bibitem{holmes2017power}
C.~C.~Holmes, S.~G.~Walker.
\newblock Assigning a value to a power likelihood in a general Bayesian model.
\newblock \emph{Biometrika}, 104(2):497--503, 2017.

\bibitem{laird1982random}
N.~M.~Laird, J.~H.~Ware.
\newblock Random-effects models for longitudinal data.
\newblock \emph{Biometrics}, 38(4):963--974, 1982.

\bibitem{breslow1993approximate}
N.~E.~Breslow, D.~G.~Clayton.
\newblock Approximate inference in generalized linear mixed models.
\newblock \emph{Journal of the American Statistical Association}, 88(421):9--25, 1993.

\bibitem{breslow1995bias}
N.~E.~Breslow, X.~Lin.
\newblock Bias correction in generalised linear mixed models with a single component of dispersion.
\newblock \emph{Biometrika}, 82(1):81--91, 1995.

\bibitem{patterson1971recovery}
H.~D.~Patterson, R.~Thompson.
\newblock Recovery of inter-block information when block sizes are unequal.
\newblock \emph{Biometrika}, 58(3):545--554, 1971.

\bibitem{kuhn2004coupling}
E.~Kuhn, M.~Lavielle.
\newblock Coupling a stochastic approximation version of EM with an MCMC procedure.
\newblock \emph{ESAIM: Probability and Statistics}, 8:115--131, 2004.

\bibitem{park2008bayesian}
T.~Park, G.~Casella.
\newblock The Bayesian lasso.
\newblock \emph{Journal of the American Statistical Association}, 103(482):681--686, 2008.

\bibitem{mundlak1978pooling}
Y.~Mundlak.
\newblock On the pooling of time series and cross section data.
\newblock \emph{Econometrica}, 46(1):69--85, 1978.

\end{thebibliography}

\end{document}
